\section{Receding-horizon model-predictive control} \label{sec:mpc}

This section closes the loop. The filter of Sections~\ref{sec:pf}--%
\ref{sec:bridge} produces a posterior over $\thetadyn$ on each
observation window; the controller of Section~\ref{sec:control}
produces a posterior over $\thetactrl$ given the current state estimate
and the dynamics-parameter posterior. The MPC driver orchestrates the
two: at each replan time, advance the plant under the previous stride's
control, refresh the filter posterior on the new observation window,
and call the controller for the next stride.

\subsection{Finite-horizon stochastic optimal control}
\label{sec:mpc-finite}

Consider a discrete-time stochastic system on a state $X_k \in \Xcal$
and a control $U_k \in \Ucal$, with transition kernel
$X_{k+1} \sim p(\,\cdot \mid X_k,\, U_k)$, a stage cost
$\ell: \Xcal \times \Ucal \to \R$, and a terminal cost
$V_f: \Xcal \to \R$. The \emph{open-loop} finite-horizon problem of
length $H$ is
\begin{equation}
\min_{u_{0:H-1} \in \Ucal^H}
   \;\E\!\left[\,\sum_{k=0}^{H-1} \ell(X_k, u_k)
                + V_f(X_H)\,\Big|\,X_0 = x\right],
\label{eq:open-loop-fh}
\end{equation}
with the expectation over the joint law of $X_{1:H}$ given the chosen
sequence $u_{0:H-1}$. Denote its optimal value by $V_H(x)$ and a
minimising sequence by $u_{0:H-1}^*(x)$. The mapping
$x \mapsto u_0^*(x)$ is the open-loop optimal first-action policy.

\subsection{The receding-horizon principle} \label{sec:mpc-receding}

In \emph{receding-horizon} or \emph{model-predictive} control,
\eqref{eq:open-loop-fh} is solved repeatedly on a sliding window. At
each replan time $t_n$, given a current state estimate $\xhat_n$:
\begin{enumerate}
  \item Solve \eqref{eq:open-loop-fh} for $u_{0:H-1}^*(\xhat_n)$;
  \item Apply only the first action (or the first stride of length
        $\Delta < H$) to the plant;
  \item Observe new data, update the state estimate to $\xhat_{n+1}$,
        return to step 1.
\end{enumerate}
The closed-loop policy is a function of $\xhat_n$ alone, so the
controller is implementable. Stability of the closed loop is delicate.
Mayne and Rawlings (2000) show that under suitable assumptions on the
horizon length and the terminal cost / terminal set, the closed-loop
state remains in a positively invariant set and converges to a
neighbourhood of a desired equilibrium. The framework relies on the
\emph{long-horizon} variant of that theory: the planning horizon $H$
is chosen long enough relative to the slowest time constant in the
dynamics that the truncation error from neglecting trajectory beyond
$H$ is small.

\subsection{Stochastic + output-feedback variants} \label{sec:mpc-variants}

Three flavours of stochastic MPC are commonly distinguished.

\begin{enumerate}
  \item[(i)] \emph{Certainty-equivalent} MPC. Replace the expectation
        in \eqref{eq:open-loop-fh} by evaluation at a point estimate
        $\hat X = \E[X \mid \xhat_n]$ and solve the resulting
        deterministic problem. Cheap, but discards uncertainty.
  \item[(ii)] \emph{Expected-cost} (stochastic) MPC. Keep the
        expectation, approximate it by Monte Carlo (sampling
        trajectories under the stochastic dynamics) or by scenario
        trees.
  \item[(iii)] \emph{Output-feedback} MPC. When $X_k$ is not directly
        observed, the controller acts on a state estimate $\xhat_n$
        produced by an external estimator. For nonlinear non-Gaussian
        systems no \emph{separation principle} holds in general; the
        joint design of estimator and controller is part of the
        modelling, not a free composition.
\end{enumerate}

The \texttt{smc2fc} pipeline is in classes (ii) and (iii)
simultaneously. The state is unobserved, so a state estimator is
required: this is the rolling-window SMC$^2$ filter of
Sections~\ref{sec:smc2}--\ref{sec:bridge}. The expectation in the
cost \eqref{eq:cost} is kept, evaluated by Monte Carlo on the SDE
under common random numbers, and minimised in $\thetactrl$-space by
the control SMC$^2$ of Section~\ref{sec:control}. The full closed-
loop policy is the composition
\[
\text{(observations)}
   \xrightarrow{\text{rolling SMC$^2$ filter}} \xhat_n
   \xrightarrow{\text{control SMC$^2$, posterior mean}} \bar\theta
   \xrightarrow{\eqref{eq:schedule-decoder}} U(\,\cdot\,; \bar\theta)
   \xrightarrow{\text{plant}} \text{(new observations)}.
\]
There is no separation principle here. The fact that this composition
works at all is an empirical claim, validated on the model classes
the framework has been tested against.

\subsection{The closed-loop pipeline --- generic algorithm}
\label{sec:mpc-pipeline}

Algorithm~\ref{alg:mpc} is the framework-level statement. Specific
applications instantiate it via the bench-driver pattern of
Section~\ref{sec:mpc-driver}.

\begin{algorithm}[H]
\caption{\texttt{smc2fc} closed-loop MPC (generic)}
\label{alg:mpc}
\begin{algorithmic}[1]
\State Initialise the plant; initialise an empty observation buffer.
\For{$n = 1, \ldots, N_w$}
  \State \textbf{Filter update.} If $n = 1$, run cold-start SMC$^2$
         (Algorithm~\ref{alg:smc2}) on the current observation window
         to produce a posterior over $\thetadyn$; otherwise warm-start
         from the previous window's posterior via the Gaussian bridge
         (Section~\ref{sec:bridge}).
  \State \textbf{State extraction.} For a small set of $\thetadyn$
         particles drawn from the posterior, run the bootstrap filter
         (Algorithm~\ref{alg:sir}) forward to the stride endpoint and
         compute the importance-weighted mean of the resulting state
         vector; call this $\xhat_n$.
  \State \textbf{Control plan.} Run the control SMC$^2$ of
         Section~\ref{sec:control}, conditioned on $\xhat_n$ and the
         posterior-mean dynamics parameters $\bar\theta_{\rm dyn}$,
         to obtain a posterior-mean schedule coefficient
         $\bar\theta_{\rm ctrl}$.
  \State \textbf{Apply.} Decode
         $U(\,\cdot\,; \bar\theta_{\rm ctrl})$ via
         \eqref{eq:schedule-decoder} and apply only its first stride
         (length $\Delta_s$) to the plant; collect the resulting
         observations and append to the buffer.
\EndFor
\end{algorithmic}
\end{algorithm}

The horizon over which the controller plans (line 5) is the full
remaining macrocycle, but only the first $\Delta_s$ bins of the plan
are committed (line 6). This is the receding-horizon truncation of
Section~\ref{sec:mpc-receding} applied to whatever model the
framework is driving.

\subsection{Driver pattern: plant, filter, controller} \label{sec:mpc-driver}

Application code instantiates Algorithm~\ref{alg:mpc} via three
artefacts.

\paragraph{(a) The plant interface, \texttt{StepwisePlant}.}
The ``plant'' is whatever object the controller acts upon to produce
new observations. It is a user-supplied class with a single mutating
method:
\begin{quote}\small
\verb|advance(stride_bins, U_schedule)|
$\;\to\;$ \texttt{(trajectory, obs\_dict)}
\end{quote}
A call to \texttt{advance} integrates the SDE for \texttt{stride\_bins}
time steps under the supplied schedule, samples the per-channel
observations, and updates the plant's internal state. Pedagogically,
\texttt{StepwisePlant} is the digital embodiment of the ``real world''
against which the controller is evaluated; mathematically, a single
call samples one transition of the joint state-observation process
under the chosen control over a stride of length $\Delta_s$. The
framework does not provide a generic \texttt{StepwisePlant}; each
user model implements its own (the FSA-v5 model's
\texttt{StepwisePlant} is one example).

\paragraph{(b) The MPC driver script.}
The receding-horizon orchestrator is a user-side script. Its
execution sequence implements Algorithm~\ref{alg:mpc} step by step:
for each window $n$ it
\begin{enumerate}
  \item runs the rolling-window filter by calling
        \texttt{run\_smc\_window\_native} (cold-start, $n=1$) or
        \texttt{run\_smc\_window\_bridge\_native} (warm-start,
        $n \ge 2$), both from
        \path{smc2fc/core/jax_native_smc.py};
  \item extracts the posterior-mean dynamics parameters and the
        smoothed latent state $\xhat_n$ at the end of the window;
  \item assembles a \texttt{ControlSpec} (cost functional, horizon
        length, current state estimate, schedule decoder);
  \item calls the control SMC$^2$ loop in
        \path{smc2fc/control/tempered_smc_loop.py}, which returns a
        posterior-mean schedule coefficient $\bar\theta_{\rm ctrl}$
        and a corresponding $U$ schedule;
  \item calls
        \verb|plant.advance(stride_bins, U_schedule)| to apply only
        the first stride of that schedule;
  \item appends the resulting observations to the data buffer and
        shifts the window.
\end{enumerate}

\paragraph{(c) The posterior-mean handoff between windows.}
The smoothed state $\xhat_n$ extracted in step 2 above is what
initialises the next window's particle filter. Consecutive windows
are therefore connected \emph{at the state level} by a single point
estimate (the importance-weighted smoothed mean) and
\emph{at the parameter level} by the Gaussian bridge of
Section~\ref{sec:bridge}. This is the practical realisation of the
warm-start that gives the rolling SMC$^2$ its computational
tractability.

\subsection{Recursive feasibility and soft barriers}
\label{sec:mpc-feasibility}

State constraints in the cost (e.g.\ a fatigue ceiling, an admissible
operating envelope) are enforced through \emph{soft-barrier} penalties
of the form $\lambda_b \cdot \max(g(X_t) - g_{\max}, 0)^2$ rather than
through hard constraints. The choice is deliberate. Under stochastic
dynamics with state-dependent diffusion and unbounded driving noise,
an exact hard constraint $g(X_t) \le g_{\max}$ cannot be guaranteed
by any controller without making the admissible-control set empty for
some achievable initial conditions; this is the standard
\emph{recursive feasibility} difficulty in stochastic MPC. A soft
penalty preserves recursive feasibility --- there is always a
feasible schedule --- and produces the constraint approximately, with
the violation rate decreasing as $\lambda_b$ is increased.

The framework's chance-constraint formulation of
Section~\ref{sec:control-cost-variants} is the more principled
alternative when the ``bad set'' is precisely defined and a violation-
rate budget is the natural specification. Soft barriers and chance
constraints are not mutually exclusive; both can be combined in the
same cost functional.

\subsection{The flat-cost-surface argument} \label{sec:mpc-robustness}

The controller commits the posterior mean $\bar\theta_{\rm ctrl}$
rather than a sample from $q_\beta$ or its full distribution. One
might worry that this discards useful information, especially when
the filter posterior over $\thetadyn$ is uncertain or biased.
Empirically the policy is nevertheless robust to filter parameter
drift, and the explanation is geometric: near the optimum, the cost
$J(\thetactrl)$ is locally flat in the directions explored by the
SMC$^2$ cloud. A first-order Taylor expansion at the posterior mean
gives, for a perturbed estimate $\bar\theta + \delta$,
\begin{equation}
J(\bar\theta + \delta)
   \;\approx\; J(\bar\theta) + \tfrac{1}{2}\delta^\top H \delta,
   \qquad H = \nabla^2 J(\bar\theta),
\end{equation}
so the additional cost from using $\bar\theta$ in place of an unknown
``oracle'' optimum scales with the Hessian. When $H$ is small in the
directions of greatest filter uncertainty, the closed-loop performance
is correspondingly insensitive to filter accuracy. This is the
mechanism by which a posterior-mean controller is reasonably robust
even when the filter is rate-limited (e.g.\ early in the rolling
schedule, before enough data has accumulated to identify the dynamics
parameters tightly).
